\begin{abstract}
We investigate whether residual stream geometry in transformers simplifies as models process more in-context examples during game-playing tasks---the hypothesized ``many-shot effect.'' We train small GPT-style transformers (208K parameters) on sequences from a 3-state hidden Markov model (\messthree), Rock-Paper-Scissors with three opponent strategies, and \ttt, then measure effective dimensionality (participation ratio), PCA concentration, and linear probe accuracy as functions of context position. Our key finding is that geometric simplification is driven primarily by \emph{state-space narrowing} in the task itself, not by a universal many-shot learning mechanism. \ttt exhibits dramatic dimensionality collapse (participation ratio: $7.85 \to 1.26$, Cohen's $d = 19.98$) because later moves are increasingly constrained. In contrast, stationary processes show approximately flat geometry across context, with participation ratio stabilizing around $2.0$--$2.3$ after the first few tokens. Unexpectedly, belief state regression quality \emph{decreases} with context position in \messthree ($R^2$: $1.0 \to 0.95$, Spearman $\rho = -0.78$, $p < 0.001$), suggesting representational drift rather than refinement. These results demonstrate that residual stream dimensionality reflects the effective entropy of the data-generating process at each position, not the accumulation of in-context evidence.
\end{abstract}
